\section{Conclusion}
\label{sec:conclusion}

This paper proposed Notebook-to-OSL as a method for studying the path from situated CNC operator knowledge to validated formal representations for Digital Twin development. The paper is centred on three linked uncertainties rather than on note capture or annotation alone.

RQ1 asks how situated operator training differs from a conventional post-session interview. The proposed design preserves the two sources separately and compares shared and source-exclusive observations, contexts, interpretations, responses, exceptions, risks, and uncertainties. This makes the value of the situated setting an empirical question rather than an assumed advantage.

RQ2 asks how operator validation changes the researcher's interpretation. The method preserves the initial annotation and records whether operator review accepts, corrects, narrows, rejects, or restricts each fragment. The resulting change history is intended to show where interpretation is reliable and where researcher inference requires correction.

RQ3 asks whether validated fragments can be transformed into traceable, semantically coherent OSL representations and what gaps remain. The transformation assessment distinguishes coherent representations from explicitly incomplete, information-losing, and unsupported mappings. It therefore treats unresolved criteria and representation failures as results rather than hiding them behind syntactically plausible output.

The next step is empirical application. A pilot should conduct the situated session and matched interview, preserve the initial and reviewed interpretations, and assess each validated fragment against explicit OSL transformation criteria. Such a study can establish case-bounded evidence about elicitation differences, interpretation change, and representation adequacy. It cannot by itself establish general superiority, runtime recommendation quality, operational correctness, safety, or manufacturing benefit.
